\scanpage{21}
\mspage{19}
% archive sheet number: 20
% hand: Grothendieck
% sheet_kind: text
% element: text | top, lines 1–3 | 4.10 Parfois
4.10. Parfois on aura à considérer un vectoriel de dim finie $E$ sur $\mathbb Q_\ell$, et une représentation continue
% element: display | upper, line 4 | u: pi to Aut(E)
\[
u:\pi\longrightarrow\operatorname{Aut}_{\mathbb Q_\ell}(E)
\tag*{(4.10.1)}
\]
% element: text | upper, lines 5–8 | Le fait à retenir
Le fait à retenir est que ($\pi$ étant compact) il existe toujours un ``réseau invariant'' par $\pi$, i.e. un sous-$\mathbb Z_\ell$-module $E_{\mathbb Z_\ell}$, \deleted{tel que} libre de type fini sur $\mathbb Z_\ell$, stable par $\pi$, tel que
% element: display | upper-middle, line 9 | E_Z_l tensor Q_l isom E
\[
E_{\mathbb Z_\ell}\otimes_{\mathbb Z_\ell}\mathbb Q_\ell\xrightarrow{\sim}E,
\]
% element: text | middle, insertion before line 10 | donc u se factorise
\deleted{\unsure{donc}}\inserted{donc $u$ se factorise \inserted{(4.10.1)}}
% element: display | middle, line 10 | u_Z_l: pi to Aut_Z_l
\[
u_{\mathbb Z_\ell}:\pi\xrightarrow{\sim}\operatorname{Aut}_{\mathbb Z_\ell}(E_{\mathbb Z_\ell})
\tag*{(4.10.2)}
\]
% element: text | middle, lines 11–15 | Il s'ensuit
Il s'ensuit \inserted{donc} que pour \unsure{tt|tout|$\forall$} sous-groupe $Q$ de $\pi$ qui est $\simeq$ une limite projective de groupes finis d'ordre \unsure{premier|premiers} à $\ell$, $u(Q)$ est un sous-groupe fini de $G=\operatorname{Aut}_{\mathbb Q_\ell}(E)$, et les cor. 4.5, et 4.6\ldots\ sont applicables.

% element: text | lower-middle, line 16 | 5. Représentations quasi-unipotentes
5. \uline{Représentations quasi-unipotentes}.

% element: text | lower, lines 17–24 | Définition 5.1
\uline{Définition} 5.1 Soit $\pi$ un groupe profini; $u:\pi\longrightarrow\operatorname{Aut}_R(E)$ une représentation linéaire de $\pi$, \deleted{\unsure{par}} \unsure{autom.} d'un vectoriel de dim finie sur \inserted{un corps} $R$. On dit que $u$ est \deleted{\unsure{quasi}} \uline{unipotente} si $\forall g\in\pi$, $u(g)$ est un automorphisme unipotent de $E$, i.e. ses valeurs propres (dans une cloture alg. de $R$) sont toutes égales à $1$, ou encore son polynôme car. $\det(\deleted{\unsure{\pi}}T\operatorname{id}-u)$ est égal à $T^n-1$. On dit que
